\section{Limitations}
\label{sec:limits}

Table~\ref{tab:assumptions} lists the assumptions the evaluated receiver
depends on. The Lean theorem applies to the receipt fields, predicate forms,
and finite domain defined in Section~\ref{sec:model}. It proves the checker
correct for that supplied domain. Completeness of the domain, parsing,
canonicalization, cryptography, clocks, stores, complexity-limit enforcement,
and the receiver hook remain outside the theorem. The differential corpus
stays within those limits and cannot establish universal implementation
equivalence.

Two valid signatures establish control of two configured keys. Organizational
independence, honest policy evaluation, uncompromised key custody, and execution
inside an attested kernel require additional deployment controls. HSM-, KMS-,
or TEE-backed keys could strengthen key custody, but we did not evaluate those
configurations.

The performance results come from a deterministic single-host setup with
SQLite receipt stores and in-memory or JSON admission stores. The experiments
exclude wide-area transport, mTLS setup, partitions, concurrent load, large
receipt stores, hardware security modules, trusted execution environments, and
crash recovery during continuation consumption. The sample counts characterize
one host and do not support operational tail-latency targets.

The \PSThreatCaseCount{} negative cases cover named attacks, not an unbounded
adversary. Six of them stop in the envelope verifier or signer and three in
pure validators, before a dispatch-capable runtime surface exists. The
dispatch-capable benchmark supplies the signed SQLite receipt and
zero-invocation observation for a unanimous policy denial and the admitted-path
cost for an allow; neither should be generalized to every internal error in
Chio.

Finally, the buyer package establishes internal consistency under its
configured keys. It cannot prove remote process integrity, prevent an
authorized kernel from lying, or compel another organization to accept the
same evidence.
